% Repository fragment for integration into the Chiral Duality manuscript.
% Assumes the theorem environments and bibliography of main.tex.
% The note is deliberately explicit about provenance: manuscript input,
% consequences derived in this note, and open frontier statements are
% separated throughout.

\chapter{The Categorical Logarithm from First Principles}
\label{chap:categorical-logarithm-frontier}

\providecommand{\Defcyc}{\mathrm{Def}_{\mathrm{cyc}}}
\providecommand{\Logsc}{\mathrm{Log}_{\mathrm{sc}}}
\providecommand{\Expsc}{\mathrm{Exp}_{\mathrm{sc}}}
\providecommand{\Philog}{\Phi_{\log}}
\providecommand{\cA}{\mathcal{A}}
\providecommand{\cB}{\mathcal{B}}
\providecommand{\cC}{\mathcal{C}}
\providecommand{\cH}{\mathcal{H}}
\providecommand{\cP}{\mathcal{P}}
\providecommand{\fg}{\mathfrak{g}}
\providecommand{\CC}{\mathbb{C}}
\providecommand{\Abar}{\overline{\mathcal{A}}}
\providecommand{\Ran}{\operatorname{Ran}}
\providecommand{\Eone}{\mathbb{E}_1}
\providecommand{\MC}{\mathrm{MC}}
\providecommand{\id}{\mathrm{id}}
\providecommand{\tr}{\mathrm{tr}}
\providecommand{\Hom}{\mathrm{Hom}}
\providecommand{\HH}{\mathrm{HH}}
\providecommand{\gr}{\mathrm{gr}}
\providecommand{\Vir}{\mathrm{Vir}}
\providecommand{\DS}{\mathrm{DS}}
\providecommand{\Com}{\mathrm{Com}}
\providecommand{\Span}{\mathrm{Span}}
\providecommand{\im}{\mathrm{im}}
\providecommand{\kerim}{\mathrm{ker}}
\providecommand{\Lie}{\mathrm{Lie}}
\providecommand{\Bin}{\mathrm{Bin}}
\providecommand{\ad}{\mathrm{ad}}
\providecommand{\End}{\mathrm{End}}
\providecommand{\op}{\mathrm{op}}
\providecommand{\Sym}{\mathrm{Sym}}
\providecommand{\rank}{\mathrm{rank}}
\providecommand{\codim}{\mathrm{codim}}
\providecommand{\Ob}{\mathrm{Ob}}
\providecommand{\Def}{\mathrm{Def}}
\providecommand{\Tot}{\mathrm{Tot}}
\providecommand{\Ahat}{\widehat{A}}
\providecommand{\ghat}{\widehat{\fg}}
\providecommand{\one}{\mathbf{1}}
\providecommand{\LamHodge}{\Lambda_{\mathrm{Hodge}}}

\begin{summary}
This chapter reorganizes the logarithmic thread running through the manuscript,
separates source statements from consequences proved here, and isolates the
first genuinely non-scalar obstruction to higher-genus deformation theory.  The
structural picture is this.

\begin{enumerate}[label=\textup{(\roman*)}]
\item The extracted repository sources explicitly define
\[
\log^{\mathrm{cat}}(\cA)=\bar B_X(\cA),
\qquad
\exp^{\mathrm{cat}}(\cC)=\Omega_X(\cC),
\]
and explicitly state that the four main theorems of the manuscript are the four
fundamental properties of a logarithm.
\item On the scalar-saturated locus, the higher-genus scalar package is governed
by a single characteristic number $\kappa(\cA)$, and the genus generating
series becomes a genuine scalar logarithm
\[
\Logsc(\cA;x)=\kappa(\cA)\,\Philog(x),
\qquad
\Philog(x)=\frac{x/2}{\sin(x/2)}-1.
\]
The exponential of this series is the scalar integrated exponential transform.
\item For any honest one-parameter level family, the Kodaira--Spencer class
$\eta$ is cohomologically central:
\[
 l_2^{\mathrm{tr}}(\eta,-)=0.
\]
Hence the first non-scalar phenomenon is not linear but quadratic.
\item The first non-scalar frontier is the primitive quadratic cone
\[
B_{\kappa,\cA}(y,y)=0,
\qquad
B_{\kappa,\cA}(u,v)=l_2^{\mathrm{tr}}(u,v)+\kappa(\cA)\,l_3^{\mathrm{tr}}(\eta,u,v),
\]
whose new ingredient is the $\eta$-Hessian
\[
H_{\eta,\cA}(u,v):=l_3^{\mathrm{tr}}(\eta,u,v).
\]
\item In the semisimple admissible highest-weight sector, the vacuum admissible
affine algebra and its Drinfeld--Sokolov reductions remain scalar-rigid.  The
first admissible family with genuinely nonzero primitive sector must therefore
lie beyond that semisimple regime.
\end{enumerate}

The note is deliberately explicit about status.  Statements imported from the
repository source are marked in the prose as manuscript input; statements proved
inside this chapter are proved here; and unresolved directions are labeled as
frontier problems.
\end{summary}

\section{Orientation, provenance, and the logarithmic principle}
\label{sec:orientation-provenance}

A logarithm is not merely a function.  It is a change of geometry.  Products are
converted into sums; local inversion appears on a convergence domain; failure of
convergence reorganizes itself into branch structure; and the residue or leading
coefficient controls the local singular behavior.  The extracted repository
sources assert that the bar construction on chiral algebras does exactly this:
it is the operation that turns OPE multiplication into additive nilpotent bar
data.  The kernel is literally logarithmic, $d\log(z_i-z_j)$, and the higher
geometry of the manuscript is organized so that the slogan ``categorical
logarithm'' is not metaphor but structure.

\begin{context}[Source-state and claim taxonomy]
\label{ctx:source-state-taxonomy}
Three distinct layers must be kept separate.
\begin{enumerate}[label=\textup{(\alph*)}]
\item \emph{Repository input.}  The extracted source files
`chapters/theory/bar\_cobar\_construction.tex`,
`chapters/theory/deformation\_theory.tex`, and
`chapters/theory/higher\_genus.tex` explicitly define the categorical
logarithm and organize the manuscript by logarithmic principles.  These same
source files also contain the stronger higher-genus statements used below.
\item \emph{Derived consequences proved here.}  The scalar logarithm,
its exponential transform, the Kodaira--Spencer centrality theorem, the
quadratic frontier theorem, the explicit $\eta$-Hessian construction, and the
admissible scalar-rigidity consequences are not stated in the extracted source
in this form; they are derived in this chapter from those source inputs.
\item \emph{Open frontier.}  The first nonsemisimple admissible family for which
$H^2_{\mathrm{cyc,prim}}$ is genuinely nonzero is not identified in the source
corpus.  The correct next target is therefore explicit computation of the
$\eta$-Hessian in such a family.
\end{enumerate}
\end{context}

\begin{remark}[Source mismatch]
\label{rem:source-mismatch}
The compiled manuscript and the extracted source do not have exactly the same
status profile.  The compiled PDF still presents parts of the full universal
Maurer--Cartan package as conjectural, while the extracted source in
`higher\_genus.tex` goes further and records an explicit minimal-model formula
for the universal Maurer--Cartan class on the scalar-saturated locus and states
that the Maurer--Cartan equation is fully resolved there.  In the present
chapter we adopt the stronger extracted-source state, but we keep this mismatch
visible because the repository should eventually reconcile the two versions.
\end{remark}

\begin{summary}[Claim-status ledger]
\label{sum:claim-ledger}
\begin{center}
\renewcommand{\arraystretch}{1.2}
\begin{tabular}{p{6.1cm}p{2.2cm}p{4.9cm}}
\textbf{Claim} & \textbf{Status} & \textbf{Role here} \\
\hline
$\log^{\mathrm{cat}}=\bar B_X$, $\exp^{\mathrm{cat}}=\Omega_X$ & source & foundation \\
Four main theorems are the four properties of a logarithm & source & conceptual spine \\
Scalar/spectral/holonomy hierarchy $\kappa\to\Delta\to\Pi$ & source & higher-genus dictionary \\
Whitehead decomposition $H^2_{\mathrm{cyc}}=\CC\eta\oplus H^2_{\mathrm{cyc,prim}}$ & source & deformation split \\
Generic rigidity $H^2_{\mathrm{cyc,prim}}=0$ under semisimplicity & source & scalar-saturated regime \\
Explicit minimal universal MC class on scalar-saturated locus & source & scalar branch \\
Scalar logarithm and scalar exponential transform & proved here & scalar packaging \\
Classification of scalar multiplicative genera factoring through $\kappa$ & proved here & uniqueness theorem \\
Kodaira--Spencer centrality $l_2^{\mathrm{tr}}(\eta,-)=0$ & proved here & kills linear frontier \\
Affine normal form along the scalar branch & proved here & structural normal form \\
Quadratic frontier and $\eta$-Hessian & proved here & first non-scalar obstruction \\
Admissible scalar rigidity in semisimple highest-weight sector & proved here & moves frontier \\
First nonsemisimple admissible family with $H^2_{\mathrm{cyc,prim}}\neq 0$ & open & next target \\
\end{tabular}
\end{center}
\end{summary}

\section{The manuscript's categorical logarithm}
\label{sec:manuscript-categorical-log}

\begin{definition}[Categorical logarithm and categorical exponential]
\label{def:categorical-log-exp-note}
Let $X$ be a smooth algebraic curve.  For an augmented $\Eone$-chiral algebra
$\cA$ on $X$, define
\[
\log^{\mathrm{cat}}(\cA):=\bar B_X(\cA),
\qquad
\exp^{\mathrm{cat}}(\cC):=\Omega_X(\cC).
\]
This is the terminology explicitly adopted in the extracted source of
`bar\_cobar\_construction.tex`.
\end{definition}

\begin{summary}[The logarithmic dictionary recorded in the source]
\label{sum:logarithmic-dictionary}
The extracted source identifies the classical logarithm and the categorical
logarithm by the following correspondences.
\begin{center}
\small
\renewcommand{\arraystretch}{1.25}
\begin{tabular}{p{5.5cm}p{5.8cm}}
\textbf{Classical logarithm} & \textbf{Categorical logarithm} \\
\hline
multiplicative-to-additive map & bar construction from OPE to bar differential \\
$\exp(\log z)=z$ on a convergence domain & $\Omega_X\bar B_X(\cA)\simeq \cA$ on the Koszul locus \\
power-series expansion & PBW spectral sequence / bar filtration \\
convergence radius & Koszul locus / $E_2$-degeneration \\
analytic continuation & coderived completion \\
branch structure & complementary Lagrangians in higher genus \\
residue / leading coefficient & modular characteristic $\kappa(\cA)$ \\
monodromy of a multivalued logarithm & higher-genus curvature and holonomy package
\end{tabular}
\end{center}
The same source explicitly states that the four main theorems of the manuscript
are, respectively, existence, invertibility on the convergence domain, branch
structure, and leading coefficient for this categorical logarithm.
\end{summary}

\begin{remark}[Why this is mathematics and not branding]
\label{rem:why-not-branding}
The slogan ``categorical logarithm'' is justified for three separate reasons.
First, the source text does not merely advertise the slogan; it records an
explicit theorem-dictionary supporting it.  Second, the bar differential is
built from a genuinely logarithmic kernel, namely $d\log(z_i-z_j)$, so the
terminology is geometrically literal.  Third, the slogan predicts and compresses
real structure: PBW as Taylor expansion, Koszulness as convergence, coderived
completion as analytic continuation, and the scalar/spectral/holonomy hierarchy
as leading coefficient, Taylor series, and monodromy.  A unifying principle of
this kind is good mathematics precisely when it produces new invariants and new
proofs; the remainder of this chapter does exactly that.
\end{remark}

\begin{summary}[Three levels of logarithmic information]
\label{sum:three-levels-logarithm}
The extracted source of `deformation\_theory.tex` explicitly organizes the
higher-genus package into three levels:
\[
\kappa \longrightarrow \Delta \longrightarrow \Pi,
\]
interpreted as scalar leading coefficient, spectral Taylor expansion, and
holonomy/monodromy of the categorical logarithm.  The present chapter works
primarily at the scalar and first non-scalar levels, but the conceptual point is
that the scalar story is the first shadow of a richer spectral and holonomy
package.
\end{summary}

\section{The scalar logarithm and the scalar integrated exponential transform}
\label{sec:scalar-log-exp}

The extracted source proves that the scalar genus package of a modular Koszul
chiral algebra $\cA$ is governed by a single number $\kappa(\cA)$ and that the
genus generating series is
\[
\sum_{g\ge 1}F_g(\cA)x^{2g}=\kappa(\cA)\left(\frac{x/2}{\sin(x/2)}-1\right).
\]
This is the precise point at which the higher-genus theory becomes a genuine
logarithmic object.

\begin{definition}[Scalar logarithm and scalar integrated exponential transform]
\label{def:scalar-log-exp}
Set
\[
\Philog(x):=\frac{x/2}{\sin(x/2)}-1 \in x^2\CC[[x^2]].
\]
For a modular Koszul chiral algebra $\cA$, define its \emph{scalar logarithm}
by
\[
\Logsc(\cA;x):=\kappa(\cA)\,\Philog(x),
\]
and define its \emph{scalar integrated exponential transform} by
\[
\Expsc(\cA;x):=\exp\bigl(\Logsc(\cA;x)\bigr)
=\exp\bigl(\kappa(\cA)\,\Philog(x)\bigr).
\]
The term ``integrated exponential transform'' is introduced here; it is not a
phrase used in the source manuscript.  What the source itself names are the
categorical logarithm and the categorical exponential.
\end{definition}

\begin{proposition}[Formal properties of the scalar logarithm; \ClaimStatusProvedHere]
\label{prop:scalar-log-formal-properties}
Assume that a class of examples is equipped with a scalar characteristic
$\kappa$ satisfying the two formal rules
\[
\kappa(\cA\star\cB)=\kappa(\cA)+\kappa(\cB),
\qquad
\kappa(\cA^!)=-\kappa(\cA),
\]
for the operation $\star$ under consideration and for Koszul duality $(-)^!$.
Then
\[
\Logsc(\cA\star\cB;x)=\Logsc(\cA;x)+\Logsc(\cB;x),
\qquad
\Expsc(\cA\star\cB;x)=\Expsc(\cA;x)\Expsc(\cB;x),
\]
and
\[
\Logsc(\cA^!;x)=-\Logsc(\cA;x),
\qquad
\Expsc(\cA^!;x)=\Expsc(\cA;x)^{-1}.
\]
\end{proposition}

\begin{proof}
Substitute the identities for $\kappa$ into the definition
$\Logsc(\cA;x)=\kappa(\cA)\Philog(x)$ and exponentiate.  No further input is
required.\end{proof}

\begin{corollary}[Self-dual anomaly neutrality; \ClaimStatusProvedHere]
\label{cor:self-dual-anomaly-neutrality}
Assume $\kappa(\cA^!)=-\kappa(\cA)$.  If $\cA\simeq \cA^!$, then
\[
\kappa(\cA)=0,
\qquad
\Logsc(\cA;x)=0,
\qquad
\Expsc(\cA;x)=1,
\qquad
F_g(\cA)=0\ \,\text{for all }g\ge 1.
\]
\end{corollary}

\begin{proof}
Self-duality gives $\kappa(\cA)=\kappa(\cA^!)=-\kappa(\cA)$, hence
$\kappa(\cA)=0$.  The remaining identities are immediate from the definition of
$\Logsc$ and the source formula for the scalar genus series.\end{proof}

\begin{theorem}[Classification of scalar multiplicative genera factoring through
$\kappa$; \ClaimStatusProvedHere]
\label{thm:classification-scalar-genera}
Let
\[
\mathcal{U}:=1+x^2\CC[[x^2]]
\]
with multiplication of formal series.  Suppose that to each object $\cA$ in a
one-scalar class one assigns an element $Z(\cA;x)\in\mathcal{U}$ such that:
\begin{enumerate}[label=\textup{(\alph*)}]
\item $Z(\cA\star\cB;x)=Z(\cA;x)Z(\cB;x)$ whenever $\kappa(\cA\star\cB)=\kappa(\cA)+\kappa(\cB)$;
\item $Z(\cA;x)$ depends on $\cA$ only through $\kappa(\cA)$;
\item the dependence on $\kappa$ is formal power-series dependence.
\end{enumerate}
Then there exists a unique series $\Psi(x)\in x^2\CC[[x^2]]$ such that
\[
Z(\cA;x)=\exp\bigl(\kappa(\cA)\Psi(x)\bigr)
\]
for every $\cA$.  If, in addition, the genus coefficients are normalized to the
source formula
\[
\sum_{g\ge 1}F_g(\cA)x^{2g}=\kappa(\cA)\Philog(x),
\]
then necessarily $\Psi(x)=\Philog(x)$ and hence
\[
Z(\cA;x)=\Expsc(\cA;x).
\]
\end{theorem}

\begin{proof}
Because $Z$ factors through $\kappa$, write $Z_z(x)$ for the value of $Z$ at an
object of scalar characteristic $z$.  Multiplicativity becomes
\[
Z_{z+w}(x)=Z_z(x)Z_w(x).
\]
The formal logarithm is defined on $1+x^2\CC[[x^2]]$, so set
\[
f_z(x):=\log Z_z(x)\in x^2\CC[[x^2]].
\]
Then
\[
f_{z+w}(x)=f_z(x)+f_w(x).
\]
Expand formally in the scalar variable:
\[
f_z(x)=\sum_{n\ge 0}c_n(x)z^n.
\]
Comparing coefficients in the additive functional equation shows that all
$c_n(x)$ vanish for $n\neq 1$, hence
\[
f_z(x)=z\Psi(x)
\]
for a unique series $\Psi(x)\in x^2\CC[[x^2]]$.  Exponentiating yields the
claimed normal form.  The final statement is immediate from the source formula
for the scalar genus generating series.\end{proof}

\begin{remark}[What has and has not been integrated]
\label{rem:what-has-been-integrated}
At scalar level the logarithmic package closes completely: the leading
coefficient $\kappa$, the scalar logarithm $\Logsc$, and the scalar exponential
transform $\Expsc$ are forced.  At non-scalar level one should resist the
temptation to exponentiate too early.  The correct first non-scalar object is
not yet a closed exponential transform but a quadratic cone in primitive
cohomology.  We now turn to that deformation-theoretic frontier.
\end{remark}

\section{Repository input from higher-genus deformation theory}
\label{sec:repository-higher-genus-input}

\begin{framework}[Whitehead decomposition, generic rigidity, and explicit scalar branch]
\label{frw:whitehead-generic-rigidity}
The extracted source of `higher\_genus.tex` records the following three inputs.
\begin{enumerate}[label=\textup{(\arabic*)}]
\item \emph{Whitehead decomposition.}  For a strongly finitely generated vertex
algebra with simple affine current algebra and finite-dimensional conformal
weight spaces, the cyclic deformation cohomology decomposes as
\[
H^2_{\mathrm{cyc}}(\cA,\cA)=\CC\cdot \eta\oplus H^2_{\mathrm{cyc,prim}}(\cA),
\]
where $\eta$ is the level class and the primitive summand consists of
$\fg$-equivariant cocycles vanishing on current-containing pairs.
\item \emph{Generic rigidity.}  If the relevant $\ghat_k$-module category is
semisimple, then $H^2_{\mathrm{cyc,prim}}(\cA)=0$, hence
$\dim H^2_{\mathrm{cyc}}(\cA,\cA)=1$ and $\cA$ is scalar-saturated.
\item \emph{Explicit scalar branch.}  On the scalar-saturated locus, the minimal
cyclic $L_\infty$-model of $\Defcyc(\cA)$ has cohomology
\[
\cH\cong \fg\oplus \CC\eta,
\qquad \deg \fg=0,
\qquad \deg\eta=2,
\]
with transferred brackets
\[
l_2^{\mathrm{tr}}=[-,-]_{\fg},
\qquad
l_3^{\mathrm{tr}}=\phi,
\qquad
l_n^{\mathrm{tr}}=0\ \text{ for }n\ge 4,
\]
and universal Maurer--Cartan element
\[
\Theta_{\cA}^{\min}=\kappa(\cA)\,\eta\otimes\Lambda,
\qquad
\Lambda:=\sum_{g\ge 1}\lambda_g.
\]
The same source explains that the Maurer--Cartan equation is then solved by the
odd parity of $s^{-1}\eta$.
\end{enumerate}
\end{framework}

\begin{framework}[Functorial stability under Drinfeld--Sokolov reduction]
\label{frw:ds-functoriality}
The extracted source also records that scalar saturation is preserved by
Drinfeld--Sokolov reduction.  In particular, if an affine algebra
$V_k(\fg)$ is scalar-saturated, then for every nilpotent $f\in\fg$ the reduced
$\mathcal{W}$-algebra $\mathcal{W}^k(\fg,f)$ is scalar-saturated as well.
\end{framework}

\begin{summary}[Scalar, spectral, and holonomy packages]
\label{sum:scalar-spectral-holonomy}
The source deformation chapter interprets the higher-genus package in three
layers:
\begin{enumerate}[label=\textup{(\roman*)}]
\item the scalar layer $\kappa$, which is the leading coefficient of the
logarithm;
\item the spectral layer $\Delta$, which is the Taylor expansion of the
logarithm;
\item the holonomy layer $\Pi$, which is the monodromy or branch ambiguity.
\end{enumerate}
The results proved below identify the first genuinely non-scalar obstruction as
an $\eta$-Hessian living exactly at the first step beyond the scalar layer.
\end{summary}

\section{Kodaira--Spencer centrality and affine normal form}
\label{sec:ks-centrality-affine-normal-form}

We now derive the first genuinely new theorem package.  The point is simple but
structural: once the level parameter is treated as an honest one-parameter
family, its Kodaira--Spencer class acts trivially on deformation cohomology.
Thus the sought-for operator $l_2^{\mathrm{tr}}(\eta,-)$ is not merely small; it
vanishes for formal reasons.

\begin{setup}[One-parameter deformation complex]
\label{setup:one-parameter-dgla}
Let $(C,\partial,[-,-])$ be a graded Lie algebra equipped with a one-parameter
family of Maurer--Cartan elements $m_k\in C^1$ depending formally on a scalar
parameter $k$.  Set
\[
d_k:=\partial+[m_k,-].
\]
Assume that the underlying graded vector space and the graded Lie bracket are
independent of $k$.  Denote by
\[
\dot m_k:=\frac{\partial m_k}{\partial k}
\]
the Kodaira--Spencer derivative and by
\[
\eta_k:=[\dot m_k]\in H^2(C,d_k)
\]
its cohomology class.
\end{setup}

\begin{theorem}[Kodaira--Spencer centrality; \ClaimStatusProvedHere]
\label{thm:ks-centrality}
In the setup above, the Kodaira--Spencer class acts trivially on cohomology:
\[
\ad_{\eta_k}=0\qquad\text{on }H^\bullet(C,d_k).
\]
Equivalently, after passing to any transferred minimal $L_\infty$-model on the
cohomology, one has
\[
l_2^{\mathrm{tr}}(\eta_k,-)=0.
\]
\end{theorem}

\begin{proof}
Differentiate the Maurer--Cartan equation
\[
\partial m_k+\frac12[m_k,m_k]=0
\]
with respect to $k$.  One obtains
\[
\partial \dot m_k+[m_k,\dot m_k]=0,
\]
that is,
\[
d_k\dot m_k=0.
\]
Thus $\eta_k=[\dot m_k]$ is a well-defined cohomology class.

Now let $x\in C$ be $d_k$-closed.  Differentiate the identity $d_kx=0$ with
respect to $k$:
\[
\frac{\partial}{\partial k}(d_kx)=0.
\]
Because the bracket is independent of $k$, this becomes
\[
d_k\!\left(\frac{\partial x}{\partial k}\right)+[\dot m_k,x]=0.
\]
Hence $[\dot m_k,x]$ is $d_k$-exact.  Therefore the cohomology class
$[\dot m_k]=\eta_k$ acts trivially on $H^\bullet(C,d_k)$.  The transferred binary
bracket on the minimal model is precisely the induced bracket on cohomology, so
$l_2^{\mathrm{tr}}(\eta_k,-)=0$.\end{proof}

\begin{lemma}[Two-$\eta$ vanishing; \ClaimStatusProvedHere]
\label{lem:two-eta-vanishing}
Let $(\cH,\{l_n\}_{n\ge 1})$ be a minimal $L_\infty$-model for a deformation
problem in which the distinguished level class $\eta\in \cH^2$ has odd
suspension $s^{-1}\eta$.  Then for every $n\ge 2$ and all $u_3,\dots,u_n\in\cH$,
\[
l_n(\eta,\eta,u_3,\dots,u_n)=0.
\]
In particular,
\[
l_n(\eta,\dots,\eta)=0\qquad(n\ge 2).
\]
\end{lemma}

\begin{proof}
In the desuspension $s^{-1}\cH$, the element $s^{-1}\eta$ has odd degree.  The
$L_\infty$-operations are graded antisymmetric on the desuspended coalgebra, so
interchanging the first two identical odd inputs multiplies the expression by
$-1$ without changing its value.  Therefore it is zero.\end{proof}

\begin{theorem}[Exact affine normal form along the scalar branch;
\ClaimStatusProvedHere]
\label{thm:affine-normal-form}
Let $(\cH,\{l_n\})$ be a minimal $L_\infty$-model with distinguished class
$\eta\in \cH^2$ satisfying Lemma~\ref{lem:two-eta-vanishing}.  For every scalar
$t\in\CC$ and every $y\in\cH$, the Maurer--Cartan series satisfies
\[
\MC(t\eta+y)=\MC(y)+t\,\Xi_\eta(y),
\]
where
\[
\Xi_\eta(y):=\sum_{m\ge 0}\frac{1}{m!}l_{m+1}(\eta,y^{\otimes m}).
\]
There are no quadratic or higher powers of $t$.
\end{theorem}

\begin{proof}
Expand the Maurer--Cartan series multilinearly:
\[
\MC(t\eta+y)=\sum_{m\ge 1}\frac{1}{m!}l_m((t\eta+y)^{\otimes m}).
\]
Every term containing at least two copies of $\eta$ vanishes by
Lemma~\ref{lem:two-eta-vanishing}.  The surviving terms are exactly the pure-$y$
part and the terms containing a single $\eta$.  Collecting them yields the
formula.\end{proof}


\begin{construction}[Scalar insertion coderivation; \ClaimStatusProvedHere]
\label{cons:scalar-insertion-coderivation}
Let $Q$ denote the coderivation on the cofree conilpotent coalgebra
$T^c(s^{-1}\cH)$ determined by the Taylor coefficients $\{l_n\}_{n\ge 1}$.  Define
a second coderivation $M_\eta$ by the Taylor coefficients
\[
(M_\eta)_n(u_1,\dots,u_n):=l_{n+1}(\eta,u_1,\dots,u_n).
\]
This is the coalgebra-level operator obtained by inserting one copy of the
level class into the $L_\infty$ structure.
\end{construction}

\begin{proposition}[Square-zero insertion differential; \ClaimStatusProvedHere]
\label{prop:square-zero-insertion}
With notation as above, the twist of $Q$ by the scalar Maurer--Cartan point
$t\eta$ is exactly
\[
Q_t=Q+tM_\eta.
\]
Consequently,
\[
[Q,M_\eta]=0,
\qquad
M_\eta^2=0.
\]
\end{proposition}

\begin{proof}
The Taylor coefficient of arity $n$ in the twist by $t\eta$ is
\[
l_n^{t\eta}(u_1,\dots,u_n)=\sum_{r\ge 0}\frac{1}{r!}l_{n+r}((t\eta)^{\otimes r},u_1,\dots,u_n).
\]
By Lemma~\ref{lem:two-eta-vanishing}, all terms with $r\ge 2$ vanish, so
\[
l_n^{t\eta}(u_1,\dots,u_n)=l_n(u_1,\dots,u_n)+t\,l_{n+1}(\eta,u_1,\dots,u_n).
\]
This is precisely the statement $Q_t=Q+tM_\eta$.  Since $t\eta$ is Maurer--Cartan
for every $t$, one has $Q_t^2=0$ for every $t$.  Expanding
\[
(Q+tM_\eta)^2=Q^2+t[Q,M_\eta]+t^2M_\eta^2
\]
and comparing coefficients of $1,t,t^2$ yields the claims.\end{proof}

\begin{corollary}[Two-step scalar spectral sequence; \ClaimStatusProvedHere]
\label{cor:two-step-scalar-sseq}
Filter the family $(T^c(s^{-1}\cH)[[t]],Q_t)$ by powers of $t$.  The associated
spectral sequence has
\[
E_1\cong H^\bullet(T^c(s^{-1}\cH),Q)\otimes \CC[[t]],
\]
its first differential is induced by $M_\eta$, and all higher differentials
vanish:
\[
d_r=0\qquad(r\ge 2).
\]
Thus the scalar-direction spectral sequence collapses at $E_2$.
\end{corollary}

\begin{proof}
The twisted differential $Q_t=Q+tM_\eta$ has exactly two filtration pieces:
$Q$, which preserves the $t$-adic filtration, and $tM_\eta$, which raises it by
one.  There is no component raising filtration by two or more.  Therefore the
only nontrivial differential beyond the $Q$-cohomology stage is the first one,
and $d_r=0$ for $r\ge 2$.\end{proof}

\begin{remark}[What affine normal form means]
\label{rem:meaning-affine-normal-form}
The theorem does not say that the scalar direction is invisible; it says that it
is \emph{affine}.  The entire dependence on the scalar coordinate is captured by
a single insertion operator $\Xi_\eta$.  By Theorem~\ref{thm:ks-centrality}, the
linear unary part of that insertion vanishes on cohomology, so the first
surviving scalar/non-scalar coupling begins quadratically in primitive
variables.
\end{remark}

\section{The first non-scalar frontier: the primitive quadratic cone}
\label{sec:primitive-quadratic-cone}

The preceding results identify the true frontier.  The sought-for linear
transport operator $l_2^{\mathrm{tr}}(\eta,-)$ vanishes.  Therefore the first
place where non-scalar behavior can occur is the quadratic interaction between
primitive directions and the scalar branch.

\begin{notation}[Primitive sector and scalar basepoint]
\label{not:primitive-sector}
Write
\[
\cP_{\cA}:=H^2_{\mathrm{cyc,prim}}(\cA)
\]
for the primitive degree-two sector from the Whitehead decomposition.  Suppress,
for notational simplicity, the inert Hodge factor $\Lambda=\sum_{g\ge 1}\lambda_g$
and write
\[
x_\kappa:=\kappa(\cA)\,\eta
\]
for the scalar Maurer--Cartan point in the minimal model.
\end{notation}

\begin{theorem}[Quadratic frontier theorem; \ClaimStatusProvedHere]
\label{thm:quadratic-frontier}
Let $\cA$ be any object for which the Whitehead decomposition holds and for
which the scalar branch $x_\kappa=\kappa(\cA)\eta$ is Maurer--Cartan.  Let
$y\in \cP_{\cA}$.  Then:
\begin{enumerate}[label=\textup{(\alph*)}]
\item every primitive direction is first-order unobstructed along the scalar
branch:
\[
l_1^{x_\kappa}(y)=0;
\]
\item the first nontrivial obstruction is quadratic:
\[
\MC(x_\kappa+y)=\frac12 B_{\kappa,\cA}(y,y)+O(y^3),
\]
where
\[
B_{\kappa,\cA}(u,v):=l_2^{\mathrm{tr}}(u,v)+\kappa(\cA)\,l_3^{\mathrm{tr}}(\eta,u,v).
\]
\end{enumerate}
Consequently, the first non-scalar branch is governed by the quadratic cone
\[
\mathcal{C}_{\kappa,\cA}:=\{y\in \cP_{\cA}\,:\,B_{\kappa,\cA}(y,y)=0\}.
\]
\end{theorem}

\begin{proof}
The twisted unary bracket is
\[
l_1^{x_\kappa}(y)=\sum_{r\ge 0}\frac{1}{r!}l_{1+r}(x_\kappa^{\otimes r},y).
\]
Because the model is minimal, the $r=0$ term vanishes.  The $r=1$ term is
\[
\kappa(\cA)\,l_2^{\mathrm{tr}}(\eta,y),
\]
which vanishes by Theorem~\ref{thm:ks-centrality}.  All terms with $r\ge 2$
contain at least two copies of $\eta$ and vanish by
Lemma~\ref{lem:two-eta-vanishing}.  This proves part~\textup{(a)}.

Now expand the Maurer--Cartan series around $x_\kappa$:
\[
\MC(x_\kappa+y)=\MC(x_\kappa)+l_1^{x_\kappa}(y)+\frac12l_2^{x_\kappa}(y,y)+O(y^3).
\]
The first term vanishes because $x_\kappa$ is Maurer--Cartan, and the second
vanishes by part~\textup{(a)}.  For the quadratic term,
\[
l_2^{x_\kappa}(u,v)=\sum_{r\ge 0}\frac{1}{r!}l_{2+r}(x_\kappa^{\otimes r},u,v).
\]
The $r=0$ term is $l_2^{\mathrm{tr}}(u,v)$, the $r=1$ term is
$\kappa(\cA)l_3^{\mathrm{tr}}(\eta,u,v)$, and all $r\ge 2$ terms vanish by
Lemma~\ref{lem:two-eta-vanishing}.  This yields the stated formula.\end{proof}

\begin{definition}[$\eta$-Hessian; \ClaimStatusProvedHere]
\label{def:eta-hessian}
The \emph{$\eta$-Hessian} of the family at $\cA$ is the trilinear term
\[
H_{\eta,\cA}(u,v):=l_3^{\mathrm{tr}}(\eta,u,v),
\qquad u,v\in \cP_{\cA}.
\]
Thus the quadratic frontier tensor is
\[
B_{\kappa,\cA}=l_2^{\mathrm{tr}}|_{\cP_{\cA}\otimes \cP_{\cA}}
+\kappa(\cA)H_{\eta,\cA}.
\]
\end{definition}

\begin{proposition}[Homotopy-transfer construction of the $\eta$-Hessian;
\ClaimStatusProvedHere]
\label{prop:eta-hessian-transfer}
Let $(C,d,[-,-])$ be a dg Lie model for the cyclic deformation complex of
$\cA$, and choose a contraction
\[
(\cP_{\cA}\oplus \CC\eta) \xrightarrow{i} C \xrightarrow{p} (\cP_{\cA}\oplus \CC\eta),
\qquad dh+hd=\id_C-ip.
\]
Then $H_{\eta,\cA}$ is represented by the standard transferred rooted-tree
formula
\begin{align*}
H_{\eta,\cA}(u,v)
&=p\bigl([h[i(\eta),i(u)],i(v)]\bigr)
 +p\bigl([h[i(u),i(v)],i(\eta)]\bigr) \\
&\qquad
 +p\bigl([h[i(v),i(\eta)],i(u)]\bigr),
\end{align*}
with the usual Koszul signs in the general graded case.  In cohomological
degree two the displayed signless formula is valid after the standard shift
convention is absorbed into the transferred brackets.
\end{proposition}

\begin{proof}
For a dg Lie algebra, the transferred ternary bracket on the minimal
$L_\infty$-model is the sum over rooted binary trees with three leaves and one
internal homotopy edge.  Evaluating that formula at the ordered triple
$(\eta,u,v)$ yields the displayed expression.  Independence of the contraction is
precisely the usual homotopy-transfer invariance of the transferred
$L_\infty$-structure.\end{proof}

\begin{lemma}[Shifted symmetry on degree-two primitives;
\ClaimStatusProvedHere]
\label{lem:shifted-symmetry-H}
On the primitive degree-two sector $\cP_{\cA}\subset H^2_{\mathrm{cyc}}(\cA,\cA)$,
both $l_2^{\mathrm{tr}}(u,v)$ and $H_{\eta,\cA}(u,v)$ are symmetric in the
ordinary sense.  Consequently, the quadratic frontier tensor
$B_{\kappa,\cA}(u,v)$ is symmetric on $\cP_{\cA}$.
\end{lemma}

\begin{proof}
After desuspension, degree-two classes become odd.  The shifted graded
antisymmetry of the transferred $L_\infty$-operations therefore becomes ordinary
symmetry when restricted to two primitive degree-two inputs.  The claim follows
for both $l_2^{\mathrm{tr}}$ and $l_3^{\mathrm{tr}}(\eta,-,-)$.\end{proof}

\begin{remark}[The frontier has moved]
\label{rem:frontier-moved}
The natural question with which we began was to compute
$l_2^{\mathrm{tr}}(\eta,-)$ and its induced quotient differential.  Theorems
\ref{thm:ks-centrality} and \ref{thm:quadratic-frontier} show that this is not
the correct frontier.  Linear transport by the level class vanishes formally.
The first genuine deformation-theoretic information lives instead in the
$\eta$-Hessian $H_{\eta,\cA}$ and the resulting quadratic cone
$\mathcal{C}_{\kappa,\cA}$.
\end{remark}

\section{Admissible affine and \texorpdfstring{$\mathcal{W}$}{W} sectors}
\label{sec:admissible-sectors}

The manuscript identifies admissible levels as the point where the generic
rigidity proof stops because semisimplicity may fail.  However, once one imports
Arakawa's semisimplicity results for the admissible highest-weight category,
one can push the rigidity argument strictly further inside that sector.

\begin{theorem}[Admissible scalar rigidity in the semisimple highest-weight
sector; \ClaimStatusProvedHere]
\label{thm:admissible-scalar-rigidity}
Let $k$ be an admissible level for a simple Lie algebra $\fg$, and let
\[
\cA=V_k(\fg)=L(k\Lambda_0)
\]
be the corresponding admissible affine vertex algebra.  Consider the cyclic
first-order deformation problem internal to the semisimple admissible
highest-weight module category of \cite{Arakawa2016RationalAdmissible}.  Then
\[
H^2_{\mathrm{cyc}}(\cA,\cA)=\CC\eta,
\qquad
H^2_{\mathrm{cyc,prim}}(\cA)=0.
\]
In particular, the semisimple admissible highest-weight sector of the vacuum
admissible affine algebra is scalar-saturated.
\end{theorem}

\begin{proof}
The source theorem on Whitehead decomposition applies at all non-critical
levels, so only the vanishing of the primitive part must be addressed.  The
Stage~2 argument in the source kills primitive cyclic deformations whenever the
relevant module category is semisimple, because the algebra object then has
vanishing Hochschild $\mathrm{Ext}^2$ in that semisimple tensor category.  For
admissible $k$, Arakawa proves that the ordinary admissible highest-weight
category is semisimple and that the vacuum algebra $V_k(\fg)$ belongs to it.
Running the source Stage~2 argument inside that semisimple admissible sector
therefore gives $H^2_{\mathrm{cyc,prim}}(\cA)=0$, and the Whitehead decomposition
reduces to $H^2_{\mathrm{cyc}}(\cA,\cA)=\CC\eta$.\end{proof}

\begin{corollary}[Drinfeld--Sokolov reductions are not the first admissible
frontier; \ClaimStatusProvedHere]
\label{cor:ds-not-first-frontier}
Let $k$ be admissible and let $f\in \fg$ be nilpotent.  In the one-parameter
cyclic deformation problem inherited from the admissible affine family, the
reduced algebra
\[
\mathcal{W}^k(\fg,f)=\DS_f(V_k(\fg))
\]
is scalar-saturated.  In particular, neither the vacuum admissible affine
algebra nor its ordinary Drinfeld--Sokolov reductions provide the first
admissible example with genuinely nonzero primitive sector.
\end{corollary}

\begin{proof}
Apply Theorem~\ref{thm:admissible-scalar-rigidity} and then invoke the source
statement that scalar saturation is preserved by Drinfeld--Sokolov reduction.\end{proof}

\begin{remark}[Where the true admissible frontier must lie]
\label{rem:true-admissible-frontier}
The first admissible family with nonzero primitive sector must lie beyond the
semisimple ordinary highest-weight regime.  Natural targets include logarithmic
extensions of admissible affine algebras, admissible conformal cosets with
Jordan blocks, relaxed highest-weight or nonordinary sectors, and nonrational
$\mathcal{W}$-algebras for which the ordinary category ceases to be semisimple.
The precise first example is an open problem.
\end{remark}

\section{How to compute the first nonsemisimple admissible \texorpdfstring{$\eta$}{eta}-Hessian}
\label{sec:procedure-compute-hessian}

Once the correct frontier is identified, the next research move becomes exact.
The problem is not to guess a slogan; it is to compute a transferred ternary
bracket in a family.

\begin{construction}[The frontier package]
\label{cons:frontier-package}
For any one-parameter family $\cA_k$ for which the Whitehead decomposition
holds, define the \emph{frontier package}
\[
\mathfrak{F}(\cA_k):=
\bigl(\cP_{\cA_k},\,\eta_k,\,\kappa(\cA_k),\,H_{\eta_k,\cA_k},\,B_{\kappa,\cA_k}\bigr),
\]
where
\[
\cP_{\cA_k}=H^2_{\mathrm{cyc,prim}}(\cA_k),
\qquad
H_{\eta_k,\cA_k}(u,v)=l_3^{\mathrm{tr}}(\eta_k,u,v),
\qquad
B_{\kappa,\cA_k}=l_2^{\mathrm{tr}}+\kappa(\cA_k)H_{\eta_k,\cA_k}.
\]
This package is the first genuinely non-scalar invariant of the categorical
logarithm.
\end{construction}

\begin{procedure}[Computing the first nonsemisimple admissible $\eta$-Hessian]
\label{proc:compute-first-hessian}
Let $\cA_k$ be a candidate nonsemisimple admissible family.
\begin{enumerate}[label=\textup{Step~\arabic*:}, leftmargin=3.5em]
\item Build a dg Lie or cyclic $L_\infty$ model $C_k$ for the deformation
complex $\Defcyc(\cA_k)$ with fixed underlying graded space and explicit
$k$-dependence in the Maurer--Cartan element $m_k$.
\item Differentiate in $k$ to obtain the closed Kodaira--Spencer cocycle
$\dot m_k$ representing the level class $\eta_k$.
\item Compute the primitive subspace
$\cP_{\cA_k}=H^2_{\mathrm{cyc,prim}}(\cA_k)$; if it vanishes, the family is not
on the frontier.
\item Choose a contraction $(p,i,h)$ from $C_k$ onto
$\cP_{\cA_k}\oplus\CC\eta_k$.
\item Evaluate the transferred ternary tree formula of
Proposition~\ref{prop:eta-hessian-transfer} to obtain
$H_{\eta_k,\cA_k}$.
\item Form the quadratic frontier tensor
$B_{\kappa,\cA_k}=l_2^{\mathrm{tr}}+\kappa(\cA_k)H_{\eta_k,\cA_k}$.
\item Solve the quadratic cone equation
\[
B_{\kappa,\cA_k}(y,y)=0
\]
inside $\cP_{\cA_k}$ to determine whether the scalar branch sprouts genuine
non-scalar branches.
\end{enumerate}
\end{procedure}

\begin{remark}[A pragmatic reading of the procedure]
\label{rem:pragmatic-reading-procedure}
The procedure is intentionally local.  One does not need the full higher-genus
nonlinear theory in order to detect the first non-scalar effect.  It is enough
to compute the primitive degree-two cohomology, the transferred ternary bracket
with one $\eta$-input, and the resulting quadratic form.  This is the exact
analogue of determining the first nontrivial Hessian of a singularity once the
linear term has been killed.
\end{remark}

\section{The platonic form of the theory}
\label{sec:platonic-form}

The preceding sections allow the entire thread to be rewritten from first
principles.  The point is not to decorate the existing manuscript with new
terminology.  The point is to reveal the invariant core to which the various
constructions, examples, and higher-genus packages already belong.

\begin{maintheorem}[Platonic form of the categorical logarithm;
\ClaimStatusProvedHere]
\label{mthm:platonic-form}
Assume the source-state recorded in
Section~\ref{sec:repository-higher-genus-input}.  Then the theory of higher-genus
bar--cobar duality organizes itself into the following hierarchy.
\begin{enumerate}[label=\textup{(\roman*)}]
\item \emph{Logarithmic linearization.}
The bar construction is the categorical logarithm and the cobar construction is
the categorical exponential.  The four main theorems are, respectively,
existence, local inversion, branch structure, and leading coefficient for this
logarithm.
\item \emph{Scalar package.}
On the scalar-saturated locus the entire scalar genus tower is governed by a
single characteristic number $\kappa(\cA)$, equivalently by the scalar
logarithm
\[
\Logsc(\cA;x)=\kappa(\cA)\left(\frac{x/2}{\sin(x/2)}-1\right).
\]
\item \emph{Vanishing of the linear frontier.}
For every honest one-parameter level family, the Kodaira--Spencer class is
cohomologically central, so
\[
l_2^{\mathrm{tr}}(\eta,-)=0.
\]
Thus there is no first-order non-scalar obstruction along the scalar branch.
\item \emph{First non-scalar invariant.}
The first genuine departure from scalarity is the $\eta$-Hessian
\[
H_{\eta,\cA}(u,v)=l_3^{\mathrm{tr}}(\eta,u,v),
\]
and the first non-scalar branch is governed by the quadratic cone
\[
B_{\kappa,\cA}(y,y)=0,
\qquad
B_{\kappa,\cA}=l_2^{\mathrm{tr}}+\kappa(\cA)H_{\eta,\cA}.
\]
\item \emph{Actual frontier.}
In semisimple admissible highest-weight sectors the primitive cohomology still
vanishes, so the first admissible family carrying genuine non-scalar data must
be nonsemisimple.  The next decisive computation is therefore explicit
construction of $H_{\eta,\cA}$ for the first such family.
\end{enumerate}
\end{maintheorem}

\begin{proof}
Part~\textup{(i)} is the explicit content of the extracted source in
`bar\_cobar\_construction.tex`.  Part~\textup{(ii)} is the scalar genus formula
from the extracted source, repackaged by Definition~\ref{def:scalar-log-exp}.
Part~\textup{(iii)} is Theorem~\ref{thm:ks-centrality}.  Part~\textup{(iv)} is
Theorem~\ref{thm:quadratic-frontier} together with
Definition~\ref{def:eta-hessian}.  Part~\textup{(v)} is
Theorem~\ref{thm:admissible-scalar-rigidity} and
Corollary~\ref{cor:ds-not-first-frontier}.\end{proof}

\begin{conclusion}
\label{conc:final-music}
What emerges is cleaner than the original thread.  The categorical logarithm is
a real organizing principle.  Its scalar part is completely rigid: one number,
one universal series, one exponential transform.  Its linear non-scalar part is
empty for formal reasons: the Kodaira--Spencer class is central.  The first true
frontier is therefore quadratic, and its new invariant is the $\eta$-Hessian.
That is where the manuscript stops, where the present chapter pushes the
mathematics forward, and where the next computation must begin.
\end{conclusion}

\begin{openproblem}[The first nonsemisimple admissible family]
\label{op:first-nonsemisimple-admissible-family}
Identify the first admissible one-parameter family $\cA_k$ for which
\[
H^2_{\mathrm{cyc,prim}}(\cA_k)\neq 0,
\]
construct a concrete deformation complex for that family, and compute the
transferred $\eta$-Hessian
\[
H_{\eta_k,\cA_k}(u,v)=l_3^{\mathrm{tr}}(\eta_k,u,v).
\]
This will determine the first genuine non-scalar quadratic cone of the
categorical logarithm in the admissible regime.
\end{openproblem}

% Suggested bibliography hooks for integration into the repository:
% \cite{BD04,LV12,GLZ22,Weibel94,EGNO,Arakawa2016RationalAdmissible,Ara12,Ara15}
